% Ad M. Brutum 1.15 (ad-brutum-01-15)
% Source: https://raw.githubusercontent.com/PerseusDL/canonical-latinLit/master/data/phi0474/phi059/phi0474.phi059.perseus-lat1.xml
% Fetched by scripts/fetch_latin.py

% Dateline (Perseus): Scr. Romae circ. med. in. Quint. a. 711 (43)

\ciceroLetterOpener{Cicero Bruto salutem}

\ciceroSection{1} Messalam habes. quibus igitur litteris tam accurate scriptis adsequi possum, subtilius ut explicem quae gerantur quaeque sint in re publica, quam tibi is exponet qui et optime omnia novit et elegantissime expedire et deferre ad te potest? cave enim existimes, Brute (quamquam non necesse est ea me ad te quae tibi nota sunt scribere; sed tamen tantam omnium laudum excellentiam non queo silentio praeterire), cave putes probitate, constantia, cura, studio rei publicae quicquam illi esse simile, ut eloquentia, qua mirabiliter excellit, vix in eo locum ad laudandum habere videatur; quamquam in hac ipsa sapientia plus apparet; ita gravi iudicio multaque arte se exercuit in verissimo genere dicendi. tanta autem industria est tantumque evigilat in studio ut non maxima ingenio, quod in eo summum est, gratia habenda videatur.

\ciceroSection{2} sed provehor amore. non enim id propositum est huic epistulae Messalam ut laudem, praesertim ad Brutum cui et virtus illius non minus quam mihi nota est et haec ipsa studia quae laudo notiora. quem cum a me dimittens graviter ferrem, hoc levabar uno, quod ad te tamquam ad alterum me proficiscens et officio fungebatur et laudem maximam sequebatur. sed haec hactenus.

\ciceroSection{3} venio nunc longo sane intervallo ad quandam epistulam, qua mihi multa tribuens unum reprehendebas quod in honoribus decernendis essem nimius et tamquam prodigus. tu hoc; alius fortasse, quod in animadversione poenaque durior, nisi forte utrumque tu. quod si ita est, utriusque rei meum iudicium studeo tibi esse notissimum †neque solum ut Solonis dictum usurpem† qui et sapientissimus fuit ex septem et legum scriptor solus ex septem. is rem publicam contineri duabus rebus dixit, praemio et poena. est scilicet utriusque rei modus sicut reliquarum et quaedam in utroque genere mediocritas. sed non tanta de re propositum est hoc loco disputare;

\ciceroSection{4} quid ego autem secutus hoc bello sim in sententiis dicendis aperire non ahenum puto. post interitum Caesaris et vestras memorabilis Idus Mart., Brute, quid ego praetermissum a vobis quantamque impendere rei publicae tempestatem dixerim non es oblitus. Magna pestis erat depulsa per vos, magna populi Romani macula deleta, vobis vero parta divina gloria, sed instrumentum regni delatum ad Lepidum et Antonium; quorum alter inconstantior, alter impurior, uterque pacem metuens, immicus otio. his ardentibus perturbandae rei publicae cupiditate quod opponi posset praesidium non habebamus; erexerat enim se civitas in retinenda libertate consentiens,

\ciceroSection{5} nos tum nimis acres, vos fortasse sapientius excessistis urbe ea quam liberaratis, Italiae sua vobis studia profitenti remisistis. itaque cum teneri urbem a parricidis viderem nec te in ea nec Cassium tuto esse posse eamque armis oppressam ab Antonio, mihi quoque ipsi esse excedendum putavi; taetrum enim spectaculum oppressa ab impiis civitas opitulandi potestate praecisa. sed animus idem qui semper infixus in patriae caritate discessum ab eius periculis ferre non potuit. itaque in medio Achaico cursu cum etesiarum diebus auster me in Italiam quasi dissuasor mei consili rettulisset, te vidi Veliae doluique vehementer. cedebas enim, Brute, cedebas, quoniam Stoici nostri negant fugere sapientis.

\ciceroSection{6} Romam ut veni, statim me obtuli Antoni sceleri atque dementiae. quem cum in me incitavissem, consilia imre coepi Brutina plane (vestri enim haec sunt propria sanguinis) rei publicae liberandae. longa sunt, quae restant, praetereunda; sunt enim de me; tantum dico, Caesarem hunc adulescentem, per quem adhuc sumus si verum fateri volumus, fluxisse ex fonte consiliorum meorum.

\ciceroSection{7} huic habiti a me honores nulli quidem, Brute, nisi debiti, nulli nisi necessarii. Vt enim primum libertatem revocare coepimus, cum se nondum ne Decimi quidem Bruti divina virtus ita commovisset ut iam id scire possemus, atque omne praesidium esset in puero qui a cervicibus nostris avertisset Antonium, quis honos ei non fuit decernendus? quamquam ego illi tum verborum laudem tribui eamque modicam, decrevi etiam imperium; quod quamquam videbatur illi aetati honorificum, tamen erat exercitum babenti necessarium. quid enim est sine imperio exercitus? statuam Philippus decrevit, celeritatem petitionis primo Servius, post maiorem etiam Servilius. nihil tum nimium videbatur.

\ciceroSection{8} sed nescio quo modo facilius in timore benigni quam in victoria grati reperiuntur. ego enim, D. Bruto liberato cum laetissimus ille civitati dies inluxisset idemque casu Bruti natalis esset, decrevi ut in fastis ad eum diem Bruti nomen adscriberetur, in eoque sum maiorum exemplum secutus qui hunc honorem mulieri Larentiae tribuerunt, cuius vos pontifices ad aram in Velabro sacrificium facere soletis. quod ego cum dabam Bruto, notam esse in fastis gratissimae victoriae sempiternam volebam. atque illo die cognovi paulo pluris in senatu malevolos esse quam gratos. per eos ipsos dies effudi, si ita vis, honores in mortuos, Hirtium et Pansam, Aquilam etiam. quod quis reprehendet, nisi qui deposito metu praeteriti periculi fuerit oblitus?

\ciceroSection{9} accedebat ad benefici memoriam gratam ratio illa quae etiam posteris esset salutaris. exstare enim volebam in crudelissimos hostis monimenta odi publici sempiterna. suspicor illud tibi minus probari quod a tuis familiaribus, optimis illis quidem viris sed in re publica rudibus, non probabatur, quod ut ovanti introire Caesari liceret decreverim. ego autem (sed erro fortasse nec tamen is sum ut mea me maxime delectent) nihil mihi videor hoc bello sensisse prudentius. cur autem ita sit aperiendum non est, ne magis videar providus fuisse quam gratus. hoc ipsum nimium; qua re alia videamus. D. Bruto decrevi honores, decrevi L. Planco. praeclara illa quidem ingenia quae gloria invitantur, sed senatus etiam sapiens qui qua quemque re putat, modo honesta, ad rem publicam iuvandam posse adduci hac utitur. at in Lepido reprehendimur; cui cum statuam in rostris statuissemus, idem illam evertimus. nos illum honore studuimus a furore revocare. vicit amentia levissimi hominis nostram prudentiam; nec tamen tantum in statuenda Lepidi statua factum est mali quantum in evertenda boni.

\ciceroSection{10} satis multa de honoribus; nunc de poena pauca dicenda sunt. intellexi enim ex tuis saepe litteris te in iis quos bello devicisti clementiam tuam velle laudari. existimo equidem nihil a te nisi sapienter; sed sceleris poenam praetermittere (id enim est quod vocatur ignoscere), etiam si in ceteris rebus tolerabile est, in hoc bello perniciosum puto. nullum enim bellum civile fuit in nostra re publica omnium quae memoria mea fuerunt, in quo bello non, utracumque pars vicisset, tamen aliqua forma esset futura rei publicae. hoc bello victores quam rem publicam simus habituri non facile adfirmarim, victis certe nulla umquam erit. dixi igitur sententias in Antonium, dixi in Lepidum severas neque tam ulciscendi causa quam ut et in praesenti sceleratos civis timore ab impugnanda patria deterrerem et in posterum documentum statuerem ne quis talem amentiam vellet imitari.

\ciceroSection{11} quamquam haec quidem sententia non magis mea fuit quam omnium. in qua videtur illud esse crudele, quod ad liberos qui nihil meruerunt poena pervenit. sed id et antiquum est et omnium civitatum, si quidem etiam Themistocli liberi eguerunt; et si iudicio damnatos eadem poena sequitur civis, qui potuimus leniores esse in hostis? quid autem queri quisquam potest de me, qui si vicisset acerbiorem se in me futurum fuisse confiteatur necesse est? habes rationem mearum sententiarum de hoc genere dumtaxat honoris et poenae; nam de ceteris rebus quid senserim quidque censuerim audisse te arbitror.

\ciceroSection{12} sed haec quidem non ita necessaria, illud valde necessarium, Brute, te in Italiam cum exercitu venire quam primum. summa est exspectatio tui. quod si Italiam attigeris, ad te concursus fiet omnium. Sive enim vicerimus, qui quidem pulcherrime viceramus nisi Lepidus perdere omnia et perire ipse cum suis concupivisset, tua nobis auctoritate opus est ad conlocandum aliquem civitatis statum; sive etiam nunc certamen reliquum est, maxima spes est cum in auctoritate tua tum in exercitus tui viribus. sed propera, per deos! scis quantum sit in temporibus, quantum in celeritate.

\ciceroSection{13} sororis tuae filiis quam diligenter consulam spero te ex matris et ex sororis litteris cogniturum. qua in causa maiorem habeo rationem tuae voluntatis quae mihi carissima est quam, ut quibusdam videor, constantiae meae. sed ego nulla in re malo quam in te amando constans et esse et videri.
